\documentclass[../../../include/open-logic-section]{subfiles}

\begin{document}

\olfileid{sth}{z}{pairs}
\olsection{جوړې}

راتلونکے بديهي اصل چې په پام کښې ئې نيسو، دا دے:

\begin{axiom}[جوړې]
د هر دوو سټونو $a, b$ دپاره سټ $\{a, b\}$ شتون لري.
\[
	\forall a \forall b \exists P \forall x (x \in P \liff (x = a \lor x = b))
\]
\end{axiom}

د تکراري تصور په کارولو سره دا اصل په لاندې ډول توجيه کولے شو.
فرض کړئ چې $a$ په پړاو $S$ کښې موجود دے او $b$ په پړاو $T$
کښې موجود دے. $M$ د $S$ او $T$ له پړاوونو څخه هغه يو وګرځوئ
چې وروسته راځي. نو ځکه چې $a$ او $b$ دواړه په پړاو $M$ کښې
موجود دي، سټ $\{a,b\}$ يوه ممکنه ټولګه ده چې له $M$ نه په هر
وروسته پړاو کښې موجوده کېدلے شي (يعنې له دواړو څخه له وروستني پړاو نه وروسته).

خو لږ تم شئ!{} ولې فرض کړو چې له $M$ نه وروسته کوم پړاوونه
\emph{شته}؟ که نۀ وي، نو زموږ توجيه ناکامېږي. نو د جوړو د اصل
د توجيه دپاره به د \olref[sth][z][story]{sec} په کيسه کښې يو بل
اصل هم زياتوو، يعنې:
\begin{enumerate}
	\item[] \stagessucc. وروستے پړاو نۀ شته.
\end{enumerate}
ايا دا اصل توجيه لري؟ د شوينفيلډ په کيسه کښې دا خبره
\emph{په ښکاره} نۀ وه بيان شوې چې وروستے پړاو نۀ شته.
خو بيا هم، که څه هم دا (په دقيقه مانا) زموږ په کيسه کښې يو
اضافي اصل دے، له هغې بنسټيزې مفکورې سره ښه سمون خوري چې
سټونه په پړاوونو کښې جوړېږي. په راتلونکي بحث کښې به ئې
يوازې ومنو. او په دې توګه د جوړو بديهي اصل هم منو.

په دې نوي اصل سره د ډېرو نورو سټونو شتون ثابتولے شو. د بېلګې په توګه:

\begin{prop}\ollabel{prop:pairsconsequences}
د هر دوو سټونو $a$ او $b$ دپاره لاندې سټونه شتون لري:
	\begin{enumerate}
		\item\ollabel{singleton} $\{a\}$
		\item\ollabel{binunion} $a \cup b$
		\item\ollabel{tuples} $\tuple{a, b}$
	\end{enumerate}
\end{prop}

\begin{proof}
\olref{singleton}. د جوړو له مخې $\{a, a\}$ شتون لري، چې د
غړو له مخې د برابرۍ په اساس هماغه $\{a\}$ دے.

\olref{binunion}. د جوړو له مخې $\{a, b\}$ شتون لري. اوس
$a \cup b = \bigcup
\{a, b\}$ د اتحاد له مخې شتون لري.

\olref{tuples}. د \olref{singleton} له مخې $\{a\}$ شتون لري.
د جوړو له مخې $\{a,
b\}$ شتون لري. اوس $\{\{a\}, \{a, b\}\} = \tuple{a, b}$
د جوړو د اصل د بيا کارولو له مخې شتون لري.
\end{proof}

\begin{prob}
وښيئ چې د هر درېو سټونو $a, b, c$ دپاره سټ $\{a, b, c\}$ شتون لري.
\end{prob}

\begin{prob}
وښيئ چې د هرې لړۍ سټونو $a_1, \ldots, a_n$ دپاره سټ $\{a_1, \ldots,
a_n\}$ شتون لري.
\end{prob}

%\begin{proof} By two applications of Pairs, $\{\{a_1, a_2\}, \{a_1,
%   a_3\}\}$ exists. By Union and Extensionality, $\bigcup \{\{a_1,
%   a_2\}, \{a_1, a_3\}\} = \{a_1, a_2, a_3\}$ exists. Repeat this
%   trick as often as necessary. \end{proof}

\end{document}
